%=====================================================================
%  Chapter 7 — The AI / Natural-Language Command Layer
%=====================================================================
\chapter{The Natural-Language Command Layer}
\label{chap:ai}

\section{Introduction}
The chapters so far built a robot that can be commanded through ROS\,2
actions and services. This final implementation chapter adds the highest
layer of the system: an interface that lets a person control the robot in
ordinary language --- ``wave your right hand'', ``take a few steps back'',
``stop!'' --- and have it act. The layer is built around a locally hosted
Large Language Model (LLM) and is the principal artificial-intelligence
contribution of the project. This chapter explains its guiding design
principle, then walks up the stack: the command vocabulary and its
server, the local LLM and tool-use, the dispatcher and its feedback loop,
the safety validator, and finally the web interface.

\section{Design Principle: the LLM Never Moves a Motor}
\label{sec:ai-principle}
The central design decision is a strict separation of concerns: the
language model is allowed to \emph{understand} and to \emph{choose}, but
never to \emph{actuate}. It does not generate joint angles or trajectories
of any kind. Instead it is given a fixed menu of high-level commands and
may only \emph{name} one of them and \emph{fill in its arguments}; a
separate, deterministic component is the only thing that turns a name into
a real robot motion. This keeps the unpredictable, generative part of the
system firmly outside the control loop: whatever the model produces, the
robot can only ever execute one of a small set of vetted, bounded motions.
Figure~\ref{fig:aiflow} shows the resulting flow, from the user down to
the robot.

\begin{figure}[ht]
  \centering
  \begin{tikzpicture}[
      font=\footnotesize, >=Latex, node distance=8mm and 10mm,
      b/.style={draw, rounded corners=2pt, align=center, thick,
                minimum height=9mm, text width=23mm},
    ]
    \node[b, fill=blue!8] (user) {User\\(web UI / voice)};
    \node[b, fill=orange!12, right=of user] (disp) {LLM dispatcher\\(Ollama tool-use)};
    \node[b, fill=red!8, right=of disp] (saf) {Safety\\validator};
    \node[b, fill=green!10, right=of saf] (srv) {Command server\\(actions / services)};
    \node[b, fill=gray!12, right=of srv] (rob) {Robot / sim};

    \draw[->] (user) -- node[above,font=\tiny]{text} (disp);
    \draw[->] (disp) -- node[above,font=\tiny]{tool+args} (saf);
    \draw[->] (saf) -- node[above,font=\tiny]{allowed} (srv);
    \draw[->] (srv) -- (rob);
    \draw[->] (rob.south) |- ($(disp.south)+(0,-7mm)$)
              -| node[pos=0.25,below,font=\tiny]{status / tilt} (disp.south);
  \end{tikzpicture}
  \caption{The natural-language command layer: the model chooses a command,
           the safety validator vets it, and only the deterministic command
           server actuates the robot. Robot status feeds back to inform the
           next decision.}
  \label{fig:aiflow}
\end{figure}

\section{The Command Vocabulary and its Server}
\label{sec:ai-vocab}
The robot exposes a fixed vocabulary of seven high-level commands:
\texttt{walk}, \texttt{stand\_still}, \texttt{sit\_down},
\texttt{wave\_hand}, \texttt{bow}, \texttt{stop} and \texttt{get\_status}.
These are hosted by a single \emph{command server} node as five ROS\,2
\textbf{actions} for the goal-driven motions (walking, standing, sitting,
waving, bowing) and two \textbf{services} for the instantaneous operations
(the emergency stop and the status query), using the interfaces defined in
\texttt{humanoid\_msgs} (Chapter~\ref{chap:software}). The server is the
authoritative gatekeeper of motion: it validates each goal's arguments
(for example, a walk distance between 0.1 and 3.0~m), enforces mutual
exclusion so that only one motion runs at a time, and checks for
cancellation every control cycle so that a stop takes effect within a
single tick. Each motion is played as a trajectory that always starts from
the robot's last commanded pose, so commands chain smoothly without
snapping. Internally, the server turns each command into the joint
trajectories of Chapter~\ref{chap:kinematics} and publishes them on the
controller topics of Chapter~\ref{chap:software}.

Crucially, this same vocabulary is described once more as a set of
\emph{tool schemas} --- one entry per command, listing its name, a
plain-language description, and its typed arguments --- which is exactly
the menu handed to the language model. Because the schemas and the server
share one source of truth, what the model is allowed to ask for can never
drift from what the robot can actually do.

\section{The Local Language Model and Tool-Use}
\label{sec:ai-llm}
The language understanding is performed by a Large Language Model run
\emph{locally} through Ollama --- the \texttt{llama3.2:3b} model, chosen to
fit within the 4\,GB of the machine's GPU. Running locally rather than
through a cloud service means the system has no external dependency, no
per-request cost, and keeps all interaction private and offline, which
suits an embedded robotics project.

The mechanism that connects the model to the robot is \textbf{tool-use}
(also called function-calling). Rather than asking the model to produce
free-form text and then trying to parse intentions out of it, the model is
given the vocabulary's tool schemas and is instructed, whenever the user
asks for something physical, to respond by \emph{calling} the single tool
that best fits and supplying its arguments as structured data. A short
system prompt fixes the model's behaviour: it adopts the identity of
``Abbes'', it must call a tool for any physical request and only chat
otherwise, it must never invent tools or arguments, and its reply to the
user is always one short, friendly, plain-language sentence with none of
the underlying machinery exposed. The model thus acts purely as a
\emph{translator} from fuzzy human language to a precise, bounded command
--- the role that language models are genuinely good at, and nothing more.

\section{The Dispatcher and the Feedback Loop}
\label{sec:ai-dispatcher}
The \emph{dispatcher} is the deterministic component that mediates between
the model and the robot. It mirrors the command server with the matching
set of action and service clients, and exposes a single uniform entry
point that takes a tool name and arguments, runs the corresponding command
to completion, and returns a uniform result (success or failure, with a
message).

Rather than a single question-and-answer exchange, the dispatcher runs a
bounded \emph{agentic loop} (Figure~\ref{fig:ailoop}). It sends the user's
message to the model; if the model calls a tool, the dispatcher executes it
on the robot and feeds the \emph{result} back into the conversation as a
tool message; the model is then queried again, and the cycle repeats until
the model stops calling tools and writes its final reply, or a fixed turn
limit is reached. The value of feeding results back is that the model can
\emph{react} to what actually happened: if a command is refused or fails,
the model sees the real reason and can recover --- for instance steadying
the robot and retrying --- rather than blindly narrating success. The
system prompt explicitly instructs the model to use these results to
recover before giving up.

\begin{figure}[ht]
  \centering
  \begin{tikzpicture}[
      font=\footnotesize, >=Latex, node distance=10mm and 14mm,
      b/.style={draw, rounded corners=2pt, align=center, thick,
                minimum height=8mm, text width=24mm},
      d/.style={draw, diamond, aspect=2, align=center, thick,
                fill=yellow!15, inner sep=1pt, text width=16mm},
    ]
    \node[b, fill=blue!8] (user) {User message};
    \node[b, fill=orange!12, right=of user] (llm) {LLM (Ollama)};
    \node[d, right=of llm] (q) {tool call?};
    \node[b, fill=green!10, below=12mm of q] (exec) {Validate \& execute on robot};
    \node[b, fill=gray!12, right=of q] (reply) {Final reply to user};

    \draw[->] (user) -- (llm);
    \draw[->] (llm) -- (q);
    \draw[->] (q) -- node[right,font=\tiny]{yes} (exec);
    \draw[->] (q) -- node[above,font=\tiny]{no} (reply);
    \draw[->] (exec) -| node[pos=0.25,below,font=\tiny]{result fed back (role: tool)} (llm.south);
  \end{tikzpicture}
  \caption{The dispatcher's bounded agentic loop: each tool result is fed
           back to the model, which can react and retry until it produces a
           final natural-language reply (subject to a turn limit).}
  \label{fig:ailoop}
\end{figure}

\section{The Safety Validator}
\label{sec:ai-safety}
Between the model's choice and its execution sits a \textbf{safety
validator} --- a thin policy gate consulted after a tool is chosen but
before it is run. If the gate denies a command, that command never becomes
a robot goal; instead the reason is fed back to the model so the robot can
\emph{explain} its refusal rather than silently doing nothing. The
validator deliberately does not duplicate the command server's per-goal
argument checks; it adds two rules that depend on physical state and on
history rather than on a single message:
\begin{itemize}[itemsep=2pt,topsep=4pt]
  \item a \textbf{tilt gate} that refuses to start \texttt{walk} when the
        robot is already tilted beyond a safe limit, reading the live tilt
        through the status query;
  \item a \textbf{rate limit} that refuses a new motion if one was allowed
        too recently, so a chatty model cannot fire a burst of rapid
        motions.
\end{itemize}
Two invariants keep the gate safe by construction: \texttt{stop} and
\texttt{get\_status} are never blocked --- the emergency brake must always
fire and a read can never be unsafe --- and the gate fails \emph{open} when
no inertial data is available, matching the command server's own
non-blocking stance. Because it is pure policy with no robot dependencies,
the validator can be unit-tested in isolation.

\section{The Web Interface}
\label{sec:ai-webui}
The whole layer is made usable through a browser-based control panel
(Figure~\ref{fig:webui}). It is served by a deliberately minimal,
dependency-free web server built on the Python standard library --- no web
framework was installed, and the root partition is tight --- which embeds
the dispatcher in-process and exposes a small set of endpoints: one to
serve the page, one to submit a command, one to poll the robot's status,
and one for the emergency stop. The page is branded around the robot's
identity, ``Abbes'', and presents an animated avatar, a live status
dashboard with a tilt indicator, clickable cards for the seven commands,
and a chat area for free-form requests.

A subtle but important design point concerns concurrency. A motion command
blocks while it runs, yet the status poll and the red \textbf{stop} button
must keep working \emph{during} a motion. The stop and status path is
therefore given its own independent ROS\,2 context, separate from the
dispatcher's, so that the emergency stop is a genuine interrupt rather than
something queued behind the running command.

\begin{figure}[ht]
  \centering
  \includegraphics[width=\textwidth]{fig7.3.png}
  \caption{The Abbes web control panel: a dependency-free browser interface
           that drives the natural-language command layer.}
  \label{fig:webui}
\end{figure}

\section{End-to-End Behaviour}
\label{sec:ai-e2e}
Putting the pieces together, a single request flows as follows. The user
types a sentence; the dispatcher sends it to the local model with the tool
menu; the model replies by calling one tool with its arguments; the safety
validator vets that choice against the robot's tilt and the recent command
history; if allowed, the command server plays the corresponding motion to
completion; the result is fed back to the model, which produces a short
spoken-style confirmation. The same path handles refusals and recovery,
and the emergency stop can interrupt any motion in progress.

\section{Conclusion}
This chapter presented the natural-language command layer, the project's
main artificial-intelligence contribution. Its design keeps the language
model strictly outside the control loop: the model only translates human
language into one of seven vetted commands, a safety validator vets that
choice, and a deterministic command server is the only component that
actuates the robot. A local LLM, an agentic feedback loop, a safety gate
and a dependency-free web interface together let a person operate the robot
conversationally while preserving the safety and determinism that a robot
demands. The next chapter evaluates the complete system --- motion, balance
and this command layer --- through simulation and testing.
